% Options for packages loaded elsewhere
\PassOptionsToPackage{unicode}{hyperref}
\PassOptionsToPackage{hyphens}{url}
%
\documentclass[
]{article}
\usepackage{lmodern}
\usepackage{amssymb,amsmath}
\usepackage{ifxetex,ifluatex}
\ifnum 0\ifxetex 1\fi\ifluatex 1\fi=0 % if pdftex
  \usepackage[T1]{fontenc}
  \usepackage[utf8]{inputenc}
  \usepackage{textcomp} % provide euro and other symbols
\else % if luatex or xetex
  \usepackage{unicode-math}
  \defaultfontfeatures{Scale=MatchLowercase}
  \defaultfontfeatures[\rmfamily]{Ligatures=TeX,Scale=1}
\fi
% Use upquote if available, for straight quotes in verbatim environments
\IfFileExists{upquote.sty}{\usepackage{upquote}}{}
\IfFileExists{microtype.sty}{% use microtype if available
  \usepackage[]{microtype}
  \UseMicrotypeSet[protrusion]{basicmath} % disable protrusion for tt fonts
}{}
\makeatletter
\@ifundefined{KOMAClassName}{% if non-KOMA class
  \IfFileExists{parskip.sty}{%
    \usepackage{parskip}
  }{% else
    \setlength{\parindent}{0pt}
    \setlength{\parskip}{6pt plus 2pt minus 1pt}}
}{% if KOMA class
  \KOMAoptions{parskip=half}}
\makeatother
\usepackage{xcolor}
\IfFileExists{xurl.sty}{\usepackage{xurl}}{} % add URL line breaks if available
\IfFileExists{bookmark.sty}{\usepackage{bookmark}}{\usepackage{hyperref}}
\hypersetup{
  pdftitle={customer},
  pdfauthor={Franklin Diaz fdiaz@paloaltonetworks.com},
  hidelinks,
  pdfcreator={LaTeX via pandoc}}
\urlstyle{same} % disable monospaced font for URLs
\usepackage{color}
\usepackage{fancyvrb}
\newcommand{\VerbBar}{|}
\newcommand{\VERB}{\Verb[commandchars=\\\{\}]}
\DefineVerbatimEnvironment{Highlighting}{Verbatim}{commandchars=\\\{\}}
% Add ',fontsize=\small' for more characters per line
\newenvironment{Shaded}{}{}
\newcommand{\AlertTok}[1]{\textcolor[rgb]{1.00,0.00,0.00}{\textbf{#1}}}
\newcommand{\AnnotationTok}[1]{\textcolor[rgb]{0.38,0.63,0.69}{\textbf{\textit{#1}}}}
\newcommand{\AttributeTok}[1]{\textcolor[rgb]{0.49,0.56,0.16}{#1}}
\newcommand{\BaseNTok}[1]{\textcolor[rgb]{0.25,0.63,0.44}{#1}}
\newcommand{\BuiltInTok}[1]{#1}
\newcommand{\CharTok}[1]{\textcolor[rgb]{0.25,0.44,0.63}{#1}}
\newcommand{\CommentTok}[1]{\textcolor[rgb]{0.38,0.63,0.69}{\textit{#1}}}
\newcommand{\CommentVarTok}[1]{\textcolor[rgb]{0.38,0.63,0.69}{\textbf{\textit{#1}}}}
\newcommand{\ConstantTok}[1]{\textcolor[rgb]{0.53,0.00,0.00}{#1}}
\newcommand{\ControlFlowTok}[1]{\textcolor[rgb]{0.00,0.44,0.13}{\textbf{#1}}}
\newcommand{\DataTypeTok}[1]{\textcolor[rgb]{0.56,0.13,0.00}{#1}}
\newcommand{\DecValTok}[1]{\textcolor[rgb]{0.25,0.63,0.44}{#1}}
\newcommand{\DocumentationTok}[1]{\textcolor[rgb]{0.73,0.13,0.13}{\textit{#1}}}
\newcommand{\ErrorTok}[1]{\textcolor[rgb]{1.00,0.00,0.00}{\textbf{#1}}}
\newcommand{\ExtensionTok}[1]{#1}
\newcommand{\FloatTok}[1]{\textcolor[rgb]{0.25,0.63,0.44}{#1}}
\newcommand{\FunctionTok}[1]{\textcolor[rgb]{0.02,0.16,0.49}{#1}}
\newcommand{\ImportTok}[1]{#1}
\newcommand{\InformationTok}[1]{\textcolor[rgb]{0.38,0.63,0.69}{\textbf{\textit{#1}}}}
\newcommand{\KeywordTok}[1]{\textcolor[rgb]{0.00,0.44,0.13}{\textbf{#1}}}
\newcommand{\NormalTok}[1]{#1}
\newcommand{\OperatorTok}[1]{\textcolor[rgb]{0.40,0.40,0.40}{#1}}
\newcommand{\OtherTok}[1]{\textcolor[rgb]{0.00,0.44,0.13}{#1}}
\newcommand{\PreprocessorTok}[1]{\textcolor[rgb]{0.74,0.48,0.00}{#1}}
\newcommand{\RegionMarkerTok}[1]{#1}
\newcommand{\SpecialCharTok}[1]{\textcolor[rgb]{0.25,0.44,0.63}{#1}}
\newcommand{\SpecialStringTok}[1]{\textcolor[rgb]{0.73,0.40,0.53}{#1}}
\newcommand{\StringTok}[1]{\textcolor[rgb]{0.25,0.44,0.63}{#1}}
\newcommand{\VariableTok}[1]{\textcolor[rgb]{0.10,0.09,0.49}{#1}}
\newcommand{\VerbatimStringTok}[1]{\textcolor[rgb]{0.25,0.44,0.63}{#1}}
\newcommand{\WarningTok}[1]{\textcolor[rgb]{0.38,0.63,0.69}{\textbf{\textit{#1}}}}
\usepackage{graphicx}
\makeatletter
\def\maxwidth{\ifdim\Gin@nat@width>\linewidth\linewidth\else\Gin@nat@width\fi}
\def\maxheight{\ifdim\Gin@nat@height>\textheight\textheight\else\Gin@nat@height\fi}
\makeatother
% Scale images if necessary, so that they will not overflow the page
% margins by default, and it is still possible to overwrite the defaults
% using explicit options in \includegraphics[width, height, ...]{}
\setkeys{Gin}{width=\maxwidth,height=\maxheight,keepaspectratio}
% Set default figure placement to htbp
\makeatletter
\def\fps@figure{htbp}
\makeatother
\setlength{\emergencystretch}{3em} % prevent overfull lines
\providecommand{\tightlist}{%
  \setlength{\itemsep}{0pt}\setlength{\parskip}{0pt}}
\setcounter{secnumdepth}{-\maxdimen} % remove section numbering
\usepackage{fancyhdr}
\pagestyle{fancy}
\fancyfoot[CO,CE]{v 0.1 | 02/15/2024 | initial version}
\fancyfoot[LE,RO]{\thepage}

\title{customer}
\author{Franklin Diaz
\href{mailto:fdiaz@paloaltonetworks.com}{\nolinkurl{fdiaz@paloaltonetworks.com}}}
\date{}

\begin{document}
\maketitle
\begin{abstract}
This repository is a set of custom tools and documentation.
\end{abstract}

\hypertarget{customer}{%
\section{Customer}\label{customer}}

\hypertarget{proxy-configuration}{%
\subsection{Proxy Configuration}\label{proxy-configuration}}

\begin{enumerate}
\def\labelenumi{\arabic{enumi}.}
\tightlist
\item
  Determine which proxy server to use, for example
  \texttt{153.2.227.107}
\item
  echo your username and pass to base64 then encrypt it, like so:
  \texttt{echo\ "username:password"\ \textbar{}\ base64\ \textbar{}\ openssl\ enc\ -aes-128-cbc\ -a}
\item
  Export this to your shell environment.
\end{enumerate}

\hypertarget{proxy-issues}{%
\subsubsection{Proxy Issues}\label{proxy-issues}}

\begin{itemize}
\tightlist
\item
  Opened case 03028240 for proxy issues in sandbox.

  \begin{itemize}
  \tightlist
  \item
    \href{https://knowledgebase.paloaltonetworks.com/KCSArticleDetail?id=kA14u0000001UtRCAU}{Configuration
    for update server when static update server is used in Firewall}
  \end{itemize}
\end{itemize}

\hypertarget{configure-bash-for-xml-api}{%
\subsection{Configure BASH for XML
API}\label{configure-bash-for-xml-api}}

Export these environment variables to prepare for script execution.

NOTE: You may also specify USER and PAN\_IP by passing flags to the
script.

\begin{Shaded}
\begin{Highlighting}[]
\BuiltInTok{export} \VariableTok{USER=}\OperatorTok{\textless{}}\ExtensionTok{the\_name\_of\_the\_firewall\_administrator\_account\_you\_configured}\OperatorTok{\textgreater{}}\NormalTok{ \# OPTIONAL}
\BuiltInTok{export} \VariableTok{PASS=}\OperatorTok{\textless{}}\ExtensionTok{the\_password\_of\_the\_firewall\_administrator\_account\_you\_configured}\OperatorTok{\textgreater{}}
\BuiltInTok{export} \VariableTok{PAN\_IP=}\OperatorTok{\textless{}}\ExtensionTok{the\_panorama\_instance\_ip\_address}\OperatorTok{\textgreater{}}\NormalTok{ \# OPTIONAL}
\end{Highlighting}
\end{Shaded}

\hypertarget{optional-bash-env-var-user}{%
\subsubsection{OPTIONAL BASH ENV VAR:
USER}\label{optional-bash-env-var-user}}

This is the username you created previously for the dedicated XML API
user.

NOTE: You may also specify a user from the command line with the
\texttt{-u} flag.

For the \texttt{USER} setting, the tool has a default value of
\texttt{xml-api-user}. To override this value, use the following command
before running the tool:

\begin{Shaded}
\begin{Highlighting}[]
\BuiltInTok{export} \VariableTok{USER=}\StringTok{"my\_new\_username"}\NormalTok{ \# }\VariableTok{set} \VariableTok{the} \VariableTok{value}
\BuiltInTok{echo} \VariableTok{$USER} \CommentTok{\# validate it was set properly}
\end{Highlighting}
\end{Shaded}

\hypertarget{bash-environment-variable-pass}{%
\subsubsection{BASH ENVIRONMENT VARIABLE:
PASS}\label{bash-environment-variable-pass}}

This is the password for the new XML API user that was created. It is
important that we do not hardcode any active passwords into the script.

For the \texttt{PASS} setting, the tool has a default value of
\texttt{default}. To override this value, use the following command
before running the tool:

\begin{Shaded}
\begin{Highlighting}[]
\BuiltInTok{export} \VariableTok{PASS=}\StringTok{"my\_secret password"}\NormalTok{ \# }\VariableTok{set} \VariableTok{the} \VariableTok{value}
\BuiltInTok{echo} \VariableTok{$PASS} \CommentTok{\# validate it was set properly}
\KeywordTok{set} \KeywordTok{|} \FunctionTok{grep} \VariableTok{$PASS} \CommentTok{\# different way to validate}
\end{Highlighting}
\end{Shaded}

Alternatively, you may choose to use a combination of the
\href{https://direnv.net/}{direnv} and
\href{https://www.passwordstore.org/}{pass} tools to make the management
of these variables slightly easier.

\hypertarget{optional-bash-env-var-pan_ip}{%
\subsubsection{OPTIONAL BASH ENV VAR:
PAN\_IP}\label{optional-bash-env-var-pan_ip}}

This is the IPv4 address of the active Panorama.

NOTE:

\begin{itemize}
\tightlist
\item
  You may specify TESTNET Panorama IP from the command line with the
  \texttt{-t} flag.
\item
  You may specify PRODUCTION Panorama IP from the command line with the
  \texttt{-p} flag.
\end{itemize}

For the \texttt{PAN\_IP} setting, the tool has a default value of
\texttt{34.134.31.136}. To override this value, use the following
command before running the tool:

\begin{Shaded}
\begin{Highlighting}[]
\BuiltInTok{export} \VariableTok{PAN\_IP=}\StringTok{"10.2.3.4"}\NormalTok{ \# }\VariableTok{example}\NormalTok{, }\VariableTok{set} \VariableTok{the} \VariableTok{value}
\BuiltInTok{echo} \VariableTok{$PAN\_IP} \CommentTok{\# validate it was set properly}
\KeywordTok{set} \KeywordTok{|} \FunctionTok{grep} \VariableTok{$PAN\_IP} \CommentTok{\# another way to validate}
\end{Highlighting}
\end{Shaded}

\hypertarget{panorama-xml-api}{%
\subsection{Panorama XML API}\label{panorama-xml-api}}

In order to allow the BASH script to communicate with Panorama via the
XML API, you need to
\href{https://docs.paloaltonetworks.com/pan-os/10-2/pan-os-panorama-api/get-started-with-the-pan-os-xml-api/enable-api-access.html}{enable
the API access}.

\textbf{As best practice, set up a separate admin account for the XML
API access.}

\begin{enumerate}
\def\labelenumi{\arabic{enumi}.}
\item
  Configure an Admin Role Profile:

  \begin{itemize}
  \tightlist
  \item
    Select \texttt{Device} \textgreater{} \texttt{Admin\ Roles} and
    click \texttt{Add}.
  \item
    Enter a \texttt{Name} to identify the role.
  \item
    Select the \texttt{Role} scope.
  \item
    In the \texttt{XML\ API} tab, click the icon for each functional
    area to toggle it to the desired setting: \texttt{Enable} or
    \texttt{Disable}.
  \item
    Click \texttt{OK} to save the profile.
  \item
    Assign the role to an administrator - complete the
    \texttt{Configure\ a\ Firewall\ Administrator\ Account} step.
  \end{itemize}
\item
  Configure a Firewall Administrator Account:

  \begin{itemize}
  \tightlist
  \item
    Select \texttt{Device} \textgreater{} \texttt{Administrators} and
    click \texttt{Add}.
  \item
    Enter a user \texttt{Name} to identify the account. If the firewall
    uses a local user database to authenticate the account, enter the
    name that you specified for the account in the database.
  \item
    Select an \texttt{Authentication\ Profile} or sequence if you
    configured either for the administrator. If the firewall uses
    \texttt{Local\ Authentication} without a local user database for the
    account, select \texttt{None} (default) and enter a
    \texttt{Password}.
  \item
    Select the \texttt{Administrator\ Type} \textgreater{}
    \texttt{Custom\ Panorama\ Admin}.
  \item
    Select the \texttt{Admin\ Role} you created for this account in the
    \texttt{Profile} section.
  \item
    Click \texttt{OK} and \texttt{Commit}.
  \end{itemize}
\end{enumerate}

NOTE: You must create the admin role and XML user in two places. The
Panorama and the FW template/template stack must both have the role and
dedicated XML user added.

\hypertarget{validate-xml-api-configuration}{%
\subsubsection{Validate XML API
configuration}\label{validate-xml-api-configuration}}

To check if the \texttt{XML\ API} has been enabled, run the following
code snippet:

\begin{Shaded}
\begin{Highlighting}[]
\ExtensionTok{curl}\NormalTok{ {-}k {-}X GET }\StringTok{\textquotesingle{}https://$PAN\_IP//api/?type=keygen\&user=$USER\&password=$PASS\textquotesingle{}}
\end{Highlighting}
\end{Shaded}

If everything was configured correctly, a response status
\texttt{success} should be returned together with an API key for the
specified account.

\hypertarget{cli-dev}{%
\subsubsection{CLI Dev}\label{cli-dev}}

\begin{itemize}
\tightlist
\item
  \href{https://docs.paloaltonetworks.com/pan-os/11-0/pan-os-panorama-api/get-started-with-the-pan-os-xml-api/explore-the-api/use-the-cli-to-find-xml-api-syntax}{Use
  the CLI to Find XML API Syntax}
\end{itemize}

\hypertarget{script-binstatus_panorama.sh}{%
\subsection{\texorpdfstring{SCRIPT:
\texttt{bin/status\_panorama.sh}}{SCRIPT: bin/status\_panorama.sh}}\label{script-binstatus_panorama.sh}}

\hypertarget{disk-status}{%
\subsubsection{Disk Status}\label{disk-status}}

The script should be updated to check disk status of log collectors.

\begin{Shaded}
\begin{Highlighting}[]
\FunctionTok{tail}\NormalTok{ follow yes  mp{-}log vld{-}1{-}0.log}
\FunctionTok{tail}\NormalTok{ follow yes  mp{-}log vld{-}2{-}0.log}
\ExtensionTok{show}\NormalTok{ system disk{-}space}
\FunctionTok{less}\NormalTok{ mp{-}log vldmgr.log}
\FunctionTok{less}\NormalTok{ mp{-}log logd.log}
\FunctionTok{less}\NormalTok{ mp{-}log configd.log }\KeywordTok{|} \ExtensionTok{match}\NormalTok{ Error}
\ExtensionTok{show}\NormalTok{ logging{-}status all}
\ExtensionTok{show}\NormalTok{ system software status }\KeywordTok{|} \ExtensionTok{match} \StringTok{"elasticsearch\textbackslash{}|es{-}"}
\ExtensionTok{show}\NormalTok{ system raid detail}
\ExtensionTok{debug}\NormalTok{ log{-}collector log{-}collection{-}stats show incoming{-}logs}
\ExtensionTok{request}\NormalTok{ log{-}fwd{-}ctrl device 010108010490 action start{-}from{-}lastack}
\end{Highlighting}
\end{Shaded}

\hypertarget{script-bindaily.sh}{%
\subsection{\texorpdfstring{SCRIPT:
\texttt{bin/daily.sh}}{SCRIPT: bin/daily.sh}}\label{script-bindaily.sh}}

\begin{itemize}
\tightlist
\item
  \href{https://knowledgebase.paloaltonetworks.com/KCSArticleDetail?id=kA10g000000ClJ9CAK}{Back
  Up Configuration and Device State from the CLI}
\end{itemize}

\hypertarget{script-bindeploy_fw_gcp.sh}{%
\subsection{\texorpdfstring{SCRIPT:
\texttt{bin/deploy\_fw\_gcp.sh}}{SCRIPT: bin/deploy\_fw\_gcp.sh}}\label{script-bindeploy_fw_gcp.sh}}

\begin{itemize}
\tightlist
\item
  Test without using a public IP. The target environment does not use
  public IP.
\item
  wait about 10 minutes after script run, login as:
  \texttt{ssh\ admin@34.67.40.130}
\end{itemize}

\hypertarget{script-binstage_release.sh}{%
\subsection{\texorpdfstring{SCRIPT:
\texttt{bin/stage\_release.sh}}{SCRIPT: bin/stage\_release.sh}}\label{script-binstage_release.sh}}

The following list can be generated by calling the script with the
\texttt{-h} flag.

\begin{Shaded}
\begin{Highlighting}[]
\ExtensionTok{franklin@ups}\NormalTok{:\textasciitilde{}/bin$ ./stage\_release.sh {-}h}

\ExtensionTok{stage\_release}\NormalTok{ script.}

\ExtensionTok{Syntax}\NormalTok{: stage\_release.sh [{-}f}\KeywordTok{|}\ExtensionTok{{-}h}\KeywordTok{|}\ExtensionTok{{-}i}\KeywordTok{|}\ExtensionTok{{-}p}\KeywordTok{|}\ExtensionTok{{-}r}\KeywordTok{|}\ExtensionTok{{-}t}\KeywordTok{|}\ExtensionTok{{-}u}\KeywordTok{|}\ExtensionTok{{-}v}\NormalTok{]}
\ExtensionTok{options}\NormalTok{:}
\ExtensionTok{{-}f}\NormalTok{     Specify a file with FW IPv4 addresses, one per line.}
\ExtensionTok{{-}h}\NormalTok{     Print this Help.}
\ExtensionTok{{-}i}\NormalTok{     Install release}
\ExtensionTok{{-}p}\NormalTok{     Use PROD environment}
\ExtensionTok{{-}r}\NormalTok{     Specify a sofwtware release to download/install (ex: 10.2.0)}
\ExtensionTok{{-}t}\NormalTok{     Use TESTNET environment}
\ExtensionTok{{-}u}\NormalTok{     Specify a USER for the Panorama}
\ExtensionTok{{-}v}\NormalTok{     Enable verbose (debug) }\ExtensionTok{mode.}
\end{Highlighting}
\end{Shaded}

\hypertarget{execution}{%
\subsubsection{Execution}\label{execution}}

Create a text file and add in IP addresses, one per line. This will
allow you to specify exactly which firewalls to operate on.

Example:

\begin{Shaded}
\begin{Highlighting}[]
\ExtensionTok{3.233.53.199}
\ExtensionTok{44.216.25.244}
\ExtensionTok{34.206.152.182}
\ExtensionTok{52.55.185.160}
\end{Highlighting}
\end{Shaded}

To run the script with a list of serials in a file, type the command
\texttt{./stage\_release.sh\ -f\ my\_ip\_list.txt}

To run the script against all firewalls that are managed by this
Panorama, type the command \texttt{stage\_release.sh}

\hypertarget{dev-notes}{%
\subsubsection{Dev Notes}\label{dev-notes}}

\begin{itemize}
\tightlist
\item
  Use \texttt{pandoc} to generate documentation
\end{itemize}

\begin{Shaded}
\begin{Highlighting}[]
\FunctionTok{sudo}\NormalTok{ apt install pandoc}
\BuiltInTok{cd}\NormalTok{ docs }\KeywordTok{\&\&} \ExtensionTok{pandoc}\NormalTok{ stage\_release.md {-}o stage\_release.pdf}
\end{Highlighting}
\end{Shaded}

\hypertarget{jump-host-setup}{%
\subsection{Jump Host Setup}\label{jump-host-setup}}

RHEL/CentOS

\begin{Shaded}
\begin{Highlighting}[]
\FunctionTok{sudo}\NormalTok{ yum install {-}y epel{-}release pandoc git}
\FunctionTok{sudo}\NormalTok{ yum install {-}y gpg2 }\CommentTok{\# for signing commits, encrypting passwords, etc.}
\FunctionTok{sudo}\NormalTok{ yum install {-}y direnv }\CommentTok{\# manage env vars}
\end{Highlighting}
\end{Shaded}

\hypertarget{pass}{%
\subsubsection{pass}\label{pass}}

the \texttt{pass} command as described in this article:
\url{https://www.redhat.com/sysadmin/management-password-store}

\hypertarget{development-environment}{%
\subsubsection{development environment}\label{development-environment}}

Run the command: sudo yum group install ``Development Tools'' Because I
would like to configure the development env as detailed here: CentOS /
RHEL 7: Install GCC (C and C++ Compiler) and Development Tools -
nixCraft (cyberciti.biz)

\hypertarget{services}{%
\subsubsection{Services}\label{services}}

\begin{itemize}
\tightlist
\item
  nginx
\item
  postfix/sendmail
\item
  crontab
\end{itemize}

\hypertarget{google-cloud}{%
\subsection{Google Cloud}\label{google-cloud}}

\begin{itemize}
\tightlist
\item
  Add detail about GCP traffic flows.
\end{itemize}

\hypertarget{gcp-production}{%
\subsubsection{GCP Production}\label{gcp-production}}

\begin{figure}
\centering
\includegraphics{../images/customer/diagram1-GCP-PROD.png}
\caption{gcp-prod}
\end{figure}

\hypertarget{gcp-testnet}{%
\subsubsection{GCP Testnet}\label{gcp-testnet}}

\begin{figure}
\centering
\includegraphics{../images/customer/diagram1-GCP-TESTNET.png}
\caption{gcp-test}
\end{figure}

\hypertarget{sso-testing-in-gcp}{%
\subsubsection{SSO Testing in GCP}\label{sso-testing-in-gcp}}

\begin{figure}
\centering
\includegraphics{../images/customer/diagram1-sandbox.png}
\caption{gcp-sand}
\end{figure}

\hypertarget{prod-and-testnet}{%
\subsection{PROD and TESTNET}\label{prod-and-testnet}}

\hypertarget{new-jersey}{%
\subsubsection{New Jersey}\label{new-jersey}}

\begin{figure}
\centering
\includegraphics{../images/customer/diagram1-New Jersey.png}
\caption{prod-newjersey}
\end{figure}

\hypertarget{georgia}{%
\subsubsection{Georgia}\label{georgia}}

\begin{figure}
\centering
\includegraphics{../images/customer/diagram1-Georgia.png}
\caption{prod-georgia}
\end{figure}

\hypertarget{github}{%
\subsection{Github}\label{github}}

\hypertarget{franklin}{%
\subsubsection{Franklin}\label{franklin}}

\hypertarget{customer-internal}{%
\subsubsection{customer internal}\label{customer-internal}}

\hypertarget{certificates}{%
\subsection{Certificates}\label{certificates}}

Ensure SSL/TLS service profile is configured under
\texttt{Setup\ \textgreater{}\ Management\ \textgreater{}\ General\ settings}.
The server certificate defined here is used to authenticate Admin users
accessing firewall management. Certificate expiration check should be
enabled too. For very strong security one should typically replace
vendor provided certificates with their certificates from their PKI
system.

\begin{itemize}
\tightlist
\item
  \href{https://github.com/PaloAltoNetworks/acme.sh}{An ACME Shell
  script: acme.sh}
\item
  \href{https://github.com/acmesh-official/acme.sh/wiki}{ACME wiki}
\end{itemize}

\begin{Shaded}
\begin{Highlighting}[]
\ExtensionTok{curl}\NormalTok{ https://raw.githubusercontent.com/acmesh{-}official/acme.sh/master/acme.sh }\KeywordTok{|} \FunctionTok{sh}\NormalTok{ {-}s {-}{-} {-}{-}install{-}online {-}m  franklin@dead10c5.org}
\ExtensionTok{show}\NormalTok{ device{-}certificate status}
\ExtensionTok{show}\NormalTok{ device{-}certificate info}
\end{Highlighting}
\end{Shaded}


\end{document}
